\noindent\textbf{Einheit 5 · Prozentuale Veränderung}

\erg{30}{a) $75\,\%$ \quad b) $90\,\%$ \quad c) $40\,\%$ \quad d) $120\,\%$ \quad e) $5\,\%$}
\erg{31}{a) $260\,€$ \quad b) $480\,€$ \quad c) 180 Mitglieder \quad d) $3{,}20\,€$ \quad e) 60 Kunden (der alte Wert steht hier hinten im Satz)}
\erg{32}{a) $220\,€$ \quad b) $45\,$kg \quad c) $60\,€$ \quad d) $81\,$m \quad e) $42\,€$ \quad f) $52\,€$ \quad g) $4{,}80\,€ \cdot 1{,}25 = 6{,}00\,€$ (nur die Erhöhung wären $1{,}20\,€$)}
\erg{33}{a) $10\,\%$ \quad b) $25\,\%$ weniger \quad c) $30\,\%$ mehr \quad d) $5\,\%$ weniger \quad e) $0{,}03 : 0{,}55 = 0{,}0545\ldots \approx 5{,}5\,\%$}
\erg{34}{a) $36\,€ \text{ sind } 90\,\%$, alter Preis $40\,€$ \quad b) $15\,€ \text{ sind } 125\,\%$, alter Preis $12\,€$ \quad c) $238\,€ : 1{,}19 = 200\,€$ \quad d) $21{,}40\,€ : 1{,}07 = 20{,}00\,€$}
\erg{35}{a) 6 Prozentpunkte \quad b) $6 : 40 = 15\,\%$ \quad c) 3 Prozentpunkte; $3 : 25 = 12\,\%$}
\erg{36}{a) $100\,$m; $8\,$m \quad b) $6 : 50 = 12\,\%$ \quad c) $21 : 250 = 0{,}084 = 8{,}4\,\%$; nein, die Rampe ist nicht zulässig}
\erg{37}{a) Ole hat die Erhöhung auf den neuen Preis bezogen; der alte Preis ist $100\,\%$, also $15 : 60 = 0{,}25 = 25\,\%$ \quad b) $9 : 45 = 20\,\%$}
\erg{38}{a) der alte Wert \quad b) nein – die Senkung wird vom höheren Zwischenwert berechnet; aus $100$ wird $110$ und daraus $99$}
